\chapter{Mozart, Wolfgang Amadeus}
\label{cha:mozart}

\href{https://en.wikipedia.org/wiki/List_of_compositions_by_Wolfgang_Amadeus_Mozart}{List of works} on Wikipedia.

\section{Sacred choral music} % (fold)
\label{sec:sacred_choral_music}

\noindent{\textbf{Masses}} \\

\noindent{Missa brevis in G (1768), K. 47d} \\ 

\noindent{Mass in C Minor, K. 427 "Great Mass" (Completed by F. Beyer)} \textit{Carolyn Sampson, Bach Collegium Japan Chorus, Bach Collegium Japan, Masaaki Suzuki}\marginnote{} \\ 

\noindent{Requiem Mass in D minor, K. 626 (1791; unfinished, completed by Franz Xaver Süssmayr)} \textit{Anne Sofie von Otter, Susan Addison, Barbara Bonney, Hans Peter Blochwitz, Willard White, English Baroque Soloists, John Eliot Gardiner}\marginnote{} \\

\noindent{\textbf{Other sacred works}} \\

\noindent{Exsultate, jubilate for Soprano in F major, K. 165/158a (1773)} \textit{Carolyn Sampson, Bach Collegium Japan, Masaaki Suzuki}\marginnote{} \\

\noindent{Vesperae solennes de Dominica in C (1779), K. 321} \textit{Mitsuko Shirai, Heidi Riess, Eberhard Büchner, Hermann Christian Polster, Rundfunkchor Leipzig, MDR-Sinfonieorchester, Herbert Kegel}\marginnote{01/2001, quality} \\

\noindent{Kyrie in D minor, K. 341/368a (1787–91)} \textit{Monteverdi Choir, English Baroque Soloists, John Eliot Gardiner}\marginnote{} \\   


% section sacred_choral_music (end)

\section{Theatre music} % (fold)
\label{sec:theatre_music}

\noindent{\textbf{Operas}} \\

\noindent{Il sogno di Scipione (1772), K. 126}  \textit{Overture: Manchester Camerata, Gábor Takács-Nagy}\marginnote{} \\
\noindent{Lucio Silla (1772), K. 135}  \textit{Overture: Manchester Camerata, Gábor Takács-Nagy}\marginnote{} \\
\noindent{La finta giardiniera (1774-1775), K. 196}  \textit{Overture: Manchester Camerata, Gábor Takács-Nagy}\marginnote{} \\
\noindent{Il re pastore (1775), K. 208}  \textit{Overture: Manchester Camerata, Gábor Takács-Nagy}\marginnote{} \\
\noindent{Zaide (1779-1780), K. 344}  \textit{Overture: Manchester Camerata, Gábor Takács-Nagy}\marginnote{} \\
\noindent{Die Entführung aus dem Serail (1781-1782), K. 384} \textit{(Arr for wind) Netherlands Wind Ensemble}\marginnote{} \\
\noindent{Der Schauspieldirektor (1786), K. 486}  \textit{Overture: Manchester Camerata, Gábor Takács-Nagy}\marginnote{} \\
\noindent{Don Giovanni, K. 527 (1787)} \textit{(Arr for wind) Netherlands Wind Ensemble}\marginnote{} \\ 

\bigskip
\noindent{\textbf{Ballet music}} \\

\bigskip
\noindent{\textbf{Arias, songs and vocal ensembles}} \\

% section theatre_music (end)

\section{Orchestral works} % (fold)
\label{sec:orchestral_works}

\noindent{\textbf{Symphonies}} \\
\href{https://en.wikipedia.org/wiki/List_of_symphonies_by_Wolfgang_Amadeus_Mozart}{List of symphonies} on Wikipedia.

\bigskip
\noindent{Symphony No. 24 in B Flat, K.182} \textit{Academy of St Martin in the Fields, Sir Neville Marriner}\marginnote{} \\
\noindent{Symphony No. 25 in G Minor, K. 183 (1773)} \textit{Academy of St Martin in the Fields, Sir Neville Marriner}\marginnote{} \\
\noindent{Symphony No. 27 in G Major, K. 199} \textit{Academy of St Martin in the Fields, Sir Neville Marriner}\marginnote{} \\
\noindent{Symphony No. 28 in C, K.200} \textit{Academy of St Martin in the Fields, Sir Neville Marriner}\marginnote{} \\
\noindent{Symphony No. 33 in B flat, K.319} \textit{Academy of St Martin in the Fields, Sir Neville Marriner}\marginnote{} \\
\noindent{Symphony No. 38 in D major "Prague", K. 504} \textit{Scottish Chamber Orchestra, Sir Charles Mackerras}\marginnote{max} \\
\noindent{Symphony No. 39 in E $\flat$ major, K. 543} \textit{Scottish Chamber Orchestra, Sir Charles Mackerras}\marginnote{max} \\
\noindent{Symphony No. 40 in G minor, K. 550} \textit{Scottish Chamber Orchestra, Sir Charles Mackerras}\marginnote{max} \\
\noindent{Symphony No. 41 in C major "Jupiter", K. 551} \textit{Scottish Chamber Orchestra, Sir Charles Mackerras}\marginnote{max} \\

\bigskip

\noindent{\textbf{Piano concertos}} \\

\noindent{Piano Concerto No. 5 in D major, K. 175 (1773)} \textit{Jean-Efflam Bavouzet, Manchester Camerata, Gábor Takács-Nagy}\marginnote{} \\
\noindent{Piano Concerto No. 6 in B$\flat$ major, K. 238} \textit{Kristian Bezuidenhout, Freiburger Barockorchester}\marginnote{} \\ 
\noindent{Piano concerto No. 7 for 3 pianos in F major "Lodron", K. 242} \textit{Karin Kei Nagano, Mari Kodama, Momo Kodama, Orchestre de la Suisse Romande, Kent Nagano}\marginnote{09/2024, max} \\
\noindent{Piano Concerto No. 8 in C major "Lützow", K. 246 (1776)} \textit{Jean-Efflam Bavouzet, Manchester Camerata, Gábor Takács-Nagy}\marginnote{} \\
\noindent{Piano concerto No. 9 in E$\flat$ major "Jenamy", K. 271} \textit{Lars Vogt, Orchestre de Chambre de Paris}\marginnote{max} \\
\noindent{Piano concerto No. 10 for 2 pianos in E$\flat$ major, K. 365} \textit{Mari Kodama, Momo Kodama, Orchestre de la Suisse Romande, Kent Nagano}\marginnote{09/2024, max} \\
\noindent{Piano Concerto No. 18 In B$\flat$ major (year), K. 456} \textit{Kristian Bezuidenhout, Freiburger Barockorchester}\marginnote{} \\ 
\noindent{Piano Concerto No. 20 in D minor, K. 466} \textit{Charles Richard-Hamelin, Les Violons du Roi, Jonathan Cohen}\marginnote{12/2023, max} \\
\noindent{Piano Concerto No. 21 in C Major, K. 467} \textit{Leif Ove Andsnes, Mahler Chamber Orchestra}\marginnote{max} \\
\noindent{Piano Concerto No. 22 in E$\flat$ major, K. 482 (1785)} \textit{Jean-Efflam Bavouzet, Manchester Camerata, Gábor Takács-Nagy}\marginnote{} \\
\noindent{Piano Concerto No. 23 in A major, K. 488} \textit{Charles Richard-Hamelin, Les Violons du Roi, Jonathan Cohen}\marginnote{12/2023, max} \\
\noindent{Piano Concerto No. 24 in C major, K. 491} \textit{Lars Vogt, Orchestre de Chambre de Paris}\marginnote{max} \\
\noindent{Piano Concerto No. 25 in C Major, K. 503}  \textit{Kristian Bezuidenhout, Freiburger Barockorchester}\marginnote{} \\
\noindent{Adagio and Fugue in C minor, K. 546}  \textit{Charles Richard-Hamelin, Les Violons du Roi, Jonathan Cohen}\marginnote{max} \\


\bigskip

\noindent{\textbf{Violin concertos}} \\

\noindent{Violin Concerto No. 1, K. 207} \textit{Gidon Kremer, Kremerata Baltica}\marginnote{high} \\
\noindent{Violin Concerto No. 2, K. 211} \textit{Gidon Kremer, Kremerata Baltica}\marginnote{high} \\
\noindent{Violin Concerto No. 3, K. 216} \textit{Gidon Kremer, Kremerata Baltica}\marginnote{high} \\
\noindent{Violin Concerto No. 4, K. 218} \textit{Gidon Kremer, Kremerata Baltica}\marginnote{high} \\
\noindent{Violin Concerto No. 5, K. 219} \textit{Gidon Kremer, Kremerata Baltica}\marginnote{high} \\

\bigskip

\noindent{\textbf{Brass concertos}} \\

\bigskip

\noindent{\textbf{Woodwind concertos}} \\

\noindent{Clarinet Concerto in A Major, K. 622} \textit{Martin Fr\"{o}st, Amsterdam Sinfonietta, Peter Oundjian}\marginnote{high} \\  

\bigskip

\noindent{\textbf{Concertante symphonies}} \\

\bigskip

\noindent{\textbf{Other concertos}} \\

% section orchestral_works (end)

\section{Piano music} % (fold)
\label{sec:piano_music}
\href{https://en.wikipedia.org/wiki/List_of_solo_piano_compositions_by_Wolfgang_Amadeus_Mozart}{List of solo piano compositions} on Wikipedia.



% section piano_music (end)

\section{Serenades, divertimenti, and other instrumental works} % (fold)
\label{sec:serenades_divertimenti_and_other_instrumental_works}

\noindent{\textbf{Serenades}} \\

\noindent{Serenade Colloredo, K. 203} \textit{Mozarteumorchester Salzburg, Roberto Gonz\'{a}lez-Monjas}\marginnote{10/2023, max} \\ 
\noindent{Serenade Notturna, K. 239} \textit{Mozarteumorchester Salzburg, Roberto Gonz\'{a}lez-Monjas}\marginnote{10/2023, max} \\ 
\noindent{Serenade No. 10 'Gran Partita', K. 361} \textit{Mark Simpson, Fraser Langton (cl), Nicholas Daniel, Emma Feilding (ob), Amy Harman, Dom Tyler (bn), Oliver Pashley, Ausiàs Garrigós Morant (basset hn), Ben Goldscheider, Angela Barnes, James Pillai, Fabian van de Geest (hn), David Stark (db)}\marginnote{high} \\  
\noindent{Serenade No. 13 for string quartet and double bass in G major, "Eine kleine Nachtmusik", K. 525 (1787)} \textit{Berliner Philharmoniker, Herbert von Karajan}\marginnote{} \\

\bigskip

\noindent{\textbf{Divertimenti}} \\

\noindent{Divertimento No. 15 in B$\flat$ major, K. 287/271h "Lodron No. 2" ("Lodronische Nachtmusik") (1777)} \textit{Berliner Philharmoniker, Herbert von Karajan}\marginnote{} \\

\bigskip

\noindent{\textbf{Marches}} \\

\noindent{March in D major, K. 237} \textit{Mozarteumorchester Salzburg, Roberto Gonz\'{a}lez-Monjas}\marginnote{10/2023, max} \\ 




% section serenades_divertimenti_and_other_instrumental_works (end)











